\chapter{Messaging Systems and Asynchronous Pipelines}
\label{ch:messaging-systems}

\begin{chaptergoalsbox}
\textbf{What you will learn:}
\begin{itemize}[leftmargin=12pt]
  \item Core messaging patterns: Point-to-Point queues vs. Publish-Subscribe (Pub/Sub) event streams.
  \item Delivery semantics: At-most-once, At-least-once, and Exactly-once processing (deduplication windows and idempotent consumers).
  \item Competing consumers, backpressure, prefetch buffers, and flow control.
  \item Poison messages, retry topologies, and Dead Letter Queues (DLQs).
  \item Idempotency keys, out-of-order message handling, and tombstone markers.
\end{itemize}

\textbf{What you should already understand:}
\begin{itemize}[leftmargin=12pt]
  \item TCP sockets, HTTP REST interfaces, and asynchronous job queues.
  \item Relational transaction concepts and data persistence from Volume I.
\end{itemize}

\textbf{Why this topic exists:}
Synchronous HTTP/gRPC RPC calls introduce tight temporal and spatial coupling: if the downstream service is down or experiencing high latency, the upstream caller blocks, saturating connection pools and cascading into catastrophic cluster-wide failure. Asynchronous messaging decouples producers and consumers in both time and space.

\textbf{Motivating practical problem:}
A payment processing pipeline uses an asynchronous queue to charge customer credit cards. During a transient database timeout, the queue re-delivers 10,000 unacknowledged messages. Because the consumer lacks idempotency guards, 4,200 customers are charged double, causing massive credit card dispute chargeback fees and merchant account suspension.
\end{chaptergoalsbox}

\section{Queues vs. Streams: Structural Trade-Offs}

Distributed message passing architectures generally fall into two distinct paradigms: message queues and event streams.

\begin{figure}[htbp]
\centering
\begin{tikzpicture}[
  scale=0.9,
  box/.style={rectangle, draw=headingblue, fill=codebg, thick, rounded corners=2pt, minimum width=2.4cm, minimum height=1.0cm, align=center, font=\sffamily\footnotesize},
  elem/.style={rectangle, draw=exampleBlue, fill=exampleBg, thick, minimum width=0.8cm, minimum height=0.6cm, font=\sffamily\tiny}
]
  % Queue diagram
  \node[anchor=west, font=\sffamily\bfseries\small, color=headingblue] at (-0.5, 2.2) {Message Queue (Destructive Consumer Ack)};
  \node[box] (p1) at (0, 1.2) {Producer};
  \node[elem] (m1) at (2.2, 1.2) {$M_4$};
  \node[elem] (m2) at (3.2, 1.2) {$M_3$};
  \node[elem] (m3) at (4.2, 1.2) {$M_2$};
  \node[box] (c1) at (6.5, 1.2) {Consumer};
  \draw[-{Stealth[scale=1.0]}, thick] (p1) -- (m1);
  \draw[-{Stealth[scale=1.0]}, thick] (m3) -- (c1);
  \node[anchor=north, font=\sffamily\tiny, color=warningRed] at (4.2, 0.7) {Deleted on ACK};

  % Log diagram
  \node[anchor=west, font=\sffamily\bfseries\small, color=headingblue] at (-0.5, -0.6) {Event Stream / Commit Log (Non-Destructive Offset Tracking)};
  \node[box] (p2) at (0, -1.6) {Producer};
  \node[elem] (e0) at (2.0, -1.6) {$0$};
  \node[elem] (e1) at (2.8, -1.6) {$1$};
  \node[elem] (e2) at (3.6, -1.6) {$2$};
  \node[elem] (e3) at (4.4, -1.6) {$3$};
  \node[elem] (e4) at (5.2, -1.6) {$4$};

  \node[box, fill=intuitionBg] (c_a) at (7.8, -1.0) {Consumer A (Off: 4)};
  \node[box, fill=tryBg] (c_b) at (7.8, -2.2) {Consumer B (Off: 2)};

  \draw[-{Stealth[scale=1.0]}, thick] (p2) -- (e0);
  \draw[-{Stealth[scale=1.0]}, thick, color=intuitionTeal] (e4) -- (c_a);
  \draw[-{Stealth[scale=1.0]}, thick, color=tryViolet] (e2) -- (c_b);
\end{tikzpicture}
\caption{Comparison of Traditional Message Queues vs. Persistent Log Streams}
\label{fig:queue-vs-stream}
\end{figure}

\begin{table}[htbp]
\centering\small
\begin{tabular}{lp{5.5cm}p{5.5cm}}
\toprule
\textbf{Dimension} & \textbf{Traditional Message Queue (RabbitMQ, SQS)} & \textbf{Event Stream / Append Log (Kafka, Pulsar)} \\
\midrule
\textbf{State Retention} & Transient. Messages are deleted immediately upon successful consumer acknowledgment. & Persistent. Immutable append-only log retained by time (retention hours) or size. \\
\textbf{Consumer Scalability} & Competing consumers pull from a single shared queue; work is distributed dynamically. & Partition-bound. Maximum parallel consumers in a group equals the number of partitions. \\
\textbf{Replayability} & Impossible. Once acknowledged, the message ceases to exist. & Native. Any consumer can reset its offset backwards to replay historical days of data. \\
\textbf{Ordering} & Total FIFO across a single queue, but broken when competing consumers fail and requeue. & Strict per-partition ordering guaranteed indefinitely. \\
\bottomrule
\end{tabular}
\caption{Architectural Comparison: Queues vs. Streams}
\label{tab:queue-vs-stream}
\end{table}

\section{Message Delivery Semantics}

In an ideal theoretical world, systems would deliver every message exactly once, with zero loss, zero duplication, and zero delay. Across real physical networks, packet losses and timeout retransmissions create distinct delivery semantics:

\subsection{At-Most-Once Delivery}
\begin{itemize}[leftmargin=18pt]
  \item \textbf{Mechanics}: The consumer receives the message, advances its offset or sends an acknowledgment to the broker \textbf{before} processing the message.
  \item \textbf{Failure Mode}: If the consumer crashes during processing, the broker believes the message was handled and never resends it. The message is permanently lost.
  \item \textbf{Use Case}: High-frequency metrics, sensor telemetry where occasional drops are preferable to blocking backpressure.
\end{itemize}

\subsection{At-Least-Once Delivery}
\begin{itemize}[leftmargin=18pt]
  \item \textbf{Mechanics}: The consumer receives the message, fully completes processing and persistence, and \textbf{then} sends an acknowledgment to the broker.
  \item \textbf{Failure Mode}: If the consumer crashes or network packets drop during the ACK transmission, the broker assumes the consumer died and re-delivers the message to another consumer. The message is processed more than once.
  \item \textbf{Production Reality}: This is the standard delivery mode of modern enterprise systems.
\end{itemize}

\subsection{The Myth and Reality of ``Exactly-Once''}

\begin{productionwarningbox}[The Network Impossibility of Physical Exactly-Once Delivery]
Physically delivering a message over an unreliable network exactly once is equivalent to solving the Two Generals' Problem, which is proven impossible. When brokers claim ``exactly-once semantics'' (EOS, such as Kafka's transactional API), they mean \textbf{effectively-once processing}:
\begin{equation}
\text{At-Least-Once Delivery} + \text{Idempotent Consumer / Deduplication} = \text{Effectively-Once Execution}
\end{equation}
If your downstream consumer executes a non-idempotent side effect (such as sending an SMS, dispatching an email, or calling a third-party non-transactional REST API), duplicates \textbf{will} trigger multiple side effects unless guarded by an application-level idempotency key.
\end{productionwarningbox}

\begin{figure}[htbp]
\centering
\begin{tikzpicture}[
  scale=0.9,
  client/.style={rectangle, draw=headingblue, fill=codebg, thick, rounded corners=2pt, minimum width=2.4cm, minimum height=1.0cm, align=center, font=\sffamily\footnotesize},
  db/.style={cylinder, draw=exampleBlue, fill=exampleBg, thick, shape border rotate=90, minimum width=1.8cm, minimum height=1.4cm, align=center, font=\sffamily\footnotesize}
]
  \node[client] (broker) at (0, 0) {Message Broker};
  \node[client] (consumer) at (5, 0) {Consumer App};
  \node[db] (db) at (10, 0) {PostgreSQL};

  \draw[-{Stealth[scale=1.0]}, thick] (broker) -- node[above, font=\sffamily\tiny] {1. Deliver(Msg, ID=X)} (consumer);
  \draw[-{Stealth[scale=1.0]}, thick, color=exampleBlue] (consumer) -- node[above, font=\sffamily\tiny] {2. Atomic Tx (Data + ID=X)} (db);
  \draw[-{Stealth[scale=1.0]}, thick, dashed, color=warningRed] (consumer) -- node[below, font=\sffamily\tiny] {3. ACK Dropped / Crash} (broker);
  \draw[-{Stealth[scale=1.0]}, thick, color=warningRed] (broker) to [bend left=40] node[above, font=\sffamily\tiny] {4. Re-deliver(Msg, ID=X)} (consumer);
  \draw[-{Stealth[scale=1.0]}, thick, color=exampleBlue] (consumer) to [bend right=40] node[below, font=\sffamily\tiny] {5. INSERT fails (Unique ID=X) $\to$ ACK} (db);
\end{tikzpicture}
\caption{Idempotent Consumer Deduplication via Database Unique Constraints}
\label{fig:idempotent-consumer}
\end{figure}

\begin{lstlisting}[language=Python, caption={Idempotent Consumer Implementation with Unique Key Deduplication}]
import psycopg2
from psycopg2 import IntegrityError

def process_order_message(conn, message: dict):
    """
    Guarantees exactly-once side-effect execution via atomic relational transaction.
    """
    message_id = message["id"]
    order_data = message["payload"]

    with conn.cursor() as cur:
        try:
            # 1. Attempt to insert message_id into processed_messages table
            cur.execute(
                "INSERT INTO processed_messages (message_id, processed_at) VALUES (%s, NOW());",
                (message_id,)
            )

            # 2. Execute business logic within the same transaction
            cur.execute(
                """
                INSERT INTO orders (order_id, customer_id, total_amount, status)
                VALUES (%s, %s, %s, %s);
                """,
                (order_data["order_id"], order_data["customer_id"], order_data["amount"], "CONFIRMED")
            )

            conn.commit()
            print(f"Successfully processed message {message_id}")

        except IntegrityError:
            # Duplicate message detected! Unique constraint on message_id violated.
            conn.rollback()
            print(f"Message {message_id} already processed. Safely ignoring and ACKing.")
\end{lstlisting}

\section{Poison Messages and Dead Letter Queues (DLQ)}

A \textbf{poison message} is a malformed or corrupted payload that causes a consumer to throw an unhandled exception every time it is read (e.g., invalid JSON syntax, null pointer, schema mismatch, or arithmetic overflow).

If a message broker automatically re-queues unacknowledged messages immediately upon consumer failure:
\begin{enumerate}[leftmargin=18pt]
  \item Consumer reads message $M$, crashes, and rejects $M$.
  \item Broker places $M$ back at the head of the queue.
  \item The same consumer (or a peer) immediately reads $M$ again and crashes.
  \item \textbf{Catastrophic Loop}: All consumers in the cluster enter an infinite crash loop, throughput drops to zero, and the CPU spikes to 100\% processing the same poison payload.
\end{enumerate}

\begin{figure}[htbp]
\centering
\begin{tikzpicture}[
  scale=0.85,
  q/.style={rectangle, draw=headingblue, fill=codebg, thick, rounded corners=2pt, minimum width=2.4cm, minimum height=1.0cm, align=center, font=\sffamily\footnotesize},
  c/.style={rectangle, draw=exampleBlue, fill=exampleBg, thick, rounded corners=2pt, minimum width=2.2cm, minimum height=1.0cm, align=center, font=\sffamily\footnotesize},
  dlq/.style={rectangle, draw=warningRed, fill=warningBg, thick, rounded corners=2pt, minimum width=2.4cm, minimum height=1.0cm, align=center, font=\sffamily\footnotesize}
]
  \node[q] (main) at (0, 0) {\textbf{Main Queue}\\Order Events};
  \node[c] (consumer) at (5, 0) {\textbf{Consumer Worker}};
  \node[q, fill=mistakeBg] (retry) at (5, -2.5) {\textbf{Retry Queue}\\TTL + Backoff};
  \node[dlq] (dlq) at (10, 0) {\textbf{Dead Letter Queue}\\\texttt{dlq.orders.failed}};

  \draw[-{Stealth[scale=1.0]}, thick] (main) -- node[above, font=\sffamily\tiny] {Deliver} (consumer);
  \draw[-{Stealth[scale=1.0]}, thick, color=mistakeAmber] (consumer) -- node[right, font=\sffamily\tiny] {Attempt $< 3$ Fail} (retry);
  \draw[-{Stealth[scale=1.0]}, thick, color=mistakeAmber] (retry) -| node[below, font=\sffamily\tiny] {TTL Expired $\to$ Re-route} (main);
  \draw[-{Stealth[scale=1.0]}, thick, color=warningRed] (consumer) -- node[above, font=\sffamily\tiny] {Attempt $\ge 3$ Exceeded} (dlq);
\end{tikzpicture}
\caption{Production Retry Pipeline with Dead Letter Queue (DLQ) Offloading}
\label{fig:dlq-pipeline}
\end{figure}

\begin{underthehoodbox}[Production Dead Letter Queue Invariants]
A robust messaging topology must adhere to three non-negotiable rules:
\begin{itemize}[leftmargin=12pt]
  \item \textbf{Bounded Retries}: Never retry infinitely. Set a hard cap (e.g., 3 to 5 attempts) tracked via an injected message header (\texttt{x-retry-count}).
  \item \textbf{Exponential Backoff with Jitter}: Never re-queue immediately. Route transient failures to delayed retry exchanges or dead-letter topics with exponential delays (e.g., $1\text{s}, 10\text{s}, 60\text{s}$) to give upstream services recovery windows.
  \item \textbf{DLQ Diagnostics \& Alerting}: When a message exhausts retries and enters the DLQ, preserve the original payload, the full stack trace, the offending worker host, and trigger a PagerDuty alert.
\end{itemize}
\end{underthehoodbox}

\section{Summary \& Key Takeaways}
\begin{itemize}[leftmargin=18pt]
  \item Message queues provide transient, destructive delivery for dynamic work distribution, whereas event streams provide immutable, ordered, replayable append logs.
  \item At-least-once delivery is the practical foundation of distributed messaging. Truly physical exactly-once delivery across networks is physically impossible due to the Two Generals' problem.
  \item Effectively-once processing requires pairing at-least-once delivery with \textbf{idempotent consumers} using unique idempotency keys and atomic database deduplication tables.
  \item Unhandled exceptions on malformed payloads create poison message infinite crash loops that can take down entire consumer clusters.
  \item Production pipelines isolate poison messages using bounded retry counters, delayed retry queues, and Dead Letter Queues (DLQs) with diagnostic telemetry.
\end{itemize}

\section{Exercises}

\subsection{Mechanics}
A payment service receives order charge events. Construct an idempotent SQL schema and an \texttt{UPSERT} query in PostgreSQL that safely handles repeated deliveries of payment authorization events without charging the customer twice.

\subsection{Debugging}
An operations team notices an alert: \emph{RabbitMQ Memory High Watermark Exceeded (Node Alarmed)}. Messages published to \texttt{order.events} are blocked. Investigation reveals that a consumer group has 10 workers, but the queue has 500,000 unacknowledged messages. The consumer configuration contains:
\begin{lstlisting}[language=Java]
channel.basicQos(0); // Prefetch count
\end{lstlisting}
Explain why RabbitMQ ran out of memory and why consumer workers stopped acknowledging messages.

\subsection{Architecture}
Design a distributed webhook notification delivery system (like Stripe's webhook infrastructure) capable of dispatching 50 million webhooks per day to arbitrary third-party customer endpoints. Address endpoint timeouts, unresponsive customer servers, exponential retry scheduling over 72 hours, and isolation so one slow customer cannot block other customers.

\section{Detailed Solutions}

\subsection{Solution to Mechanics Exercise}
\begin{lstlisting}[language=SQL]
-- Idempotent Ledger Schema
CREATE TABLE payment_authorizations (
    idempotency_key VARCHAR(64) PRIMARY KEY,
    customer_id VARCHAR(64) NOT NULL,
    amount_cents BIGINT NOT NULL,
    status VARCHAR(32) NOT NULL,
    created_at TIMESTAMPTZ NOT NULL DEFAULT NOW()
);

-- Safe Upsert Query: returns existing record on conflict without re-charging
INSERT INTO payment_authorizations (idempotency_key, customer_id, amount_cents, status)
VALUES ('pay_tok_948fbc', 'cust_123', 4999, 'AUTHORIZED')
ON CONFLICT (idempotency_key) DO UPDATE
SET status = payment_authorizations.status
RETURNING *;
\end{lstlisting}

\subsection{Solution to Debugging Exercise}
\textbf{The basicQos(0) Memory Trap:}
In AMQP / RabbitMQ, setting \texttt{basicQos(0)} specifies an \textbf{unlimited prefetch count}.
When workers connect and invoke \texttt{basicConsume}, RabbitMQ attempts to push \emph{every single message in the queue} across the TCP socket into the client workers' in-memory buffers as fast as the network allows.
Because the workers take time to process each message, hundreds of thousands of messages are pulled out of disk/broker management into unacknowledged RAM buffers.
The broker must track state metadata for all 500,000 in-flight unacknowledged deliveries in RAM. Once broker memory crosses the 40\% watermark, RabbitMQ triggers an alarm, suspends socket reading, and stops accepting new ACKs and publishes.
\textbf{The Fix:} Always set a bounded prefetch count:
\begin{lstlisting}[language=Java]
channel.basicQos(50); // Fetch at most 50 unacked messages per consumer thread
\end{lstlisting}

\subsection{Solution to Architecture Exercise}
\textbf{Distributed Webhook Dispatch Architecture:}
\begin{enumerate}[leftmargin=18pt]
  \item \textbf{Partitioning and Customer Isolation}: Route outgoing webhooks into tiered Kafka topics or partitioned queues based on customer hash: \texttt{webhooks.fast} and \texttt{webhooks.slow}. If customer $X$'s endpoint times out continuously, their events are dynamically downgraded to the slow tier, preventing thread starvation for healthy customers.
  \item \textbf{Tiered Delayed Retry Queues}:
  \begin{itemize}[leftmargin=12pt]
    \item Retry 1: 15 seconds delay.
    \item Retry 2: 5 minutes delay.
    \item Retry 3: 1 hour delay.
    \item Retry 4--10: Bounded exponential backoff up to 72 hours.
  \end{itemize}
  Implement delays using message timestamps and separate delayed topics (or Redis sorted sets where score is unix execution epoch).
  \item \textbf{HTTP Client Hard Limits}: Set strict timeouts on outbound connections: connect timeout $3\text{ s}$, read timeout $10\text{ s}$. Enforce circuit breakers per domain: if a customer domain fails 50 consecutive requests, short-circuit immediately.
\end{enumerate}

\section{Further Reading \& Citations}
\begin{itemize}[leftmargin=18pt]
  \item Kleppmann, M. (2017). \emph{Designing Data-Intensive Applications}. O'Reilly Media. Chapter 11: Stream Processing.
  \item Hohpe, G., \& Woolf, B. (2003). \emph{Enterprise Integration Patterns}. Addison-Wesley.
  \item Akidau, T., et al. (2015). \emph{The dataflow model: a practical approach to balancing correctness, latency, and cost in massive-scale, unbounded, out-of-order data processing}. Proceedings of the VLDB Endowment, 8(12), 1792--1803.
\end{itemize}
